\documentclass[11pt,oneside]{memoir}
\usepackage[T1]{fontenc}
\usepackage{lmodern}
\chapterstyle{bianchi}
\setlrmarginsandblock{3cm}{3cm}{*}
\setulmarginsandblock{3cm}{3cm}{*}
\checkandfixthelayout

\begin{document}

\chapter{The Lighthouse Keeper}

\epigraph{The sea, once it casts its spell, holds one in its net of wonder forever.}{Jacques Cousteau}

Elias had kept the light at Carrick Point for twenty-three years, and in all that time he
had never once let it go dark. He wound the clockwork every four hours, trimmed the wick,
polished the great lens until it threw his own face back at him a hundred times over.

The storm came in on the third night of November. By dusk the waves were breaking over
the lower gallery, and the wind had a voice he had not heard before, high and thin, like a
kettle left too long on the stove.

\section{The Ship}

At a little past midnight he saw it: a single green light, low on the water to the
south-west, rising and falling with the swell. A ship, then, and close---far too close to
the Needles.

He lit the signal rockets one after another. The first was swallowed by the rain. The
second burst red above the rocks, and for a moment he saw the vessel clearly, a coastal
steamer heeling hard to port.

\section{Morning}

When the grey light came, the steamer was gone. Not wrecked, he would learn later, but
safe in the lee of Gull Island, its captain having seen the rocket and turned away just in
time.

Elias logged the incident in his neat, small hand:

\begin{quote}
\textit{3 Nov.\ 00:20. Vessel sighted SW, standing into danger. Two rockets fired.
Vessel altered course. Wind NW force 10. Light maintained throughout.}
\end{quote}

Then he climbed the hundred and twelve steps, and wound the clock again.

\end{document}
